\documentclass[11pt]{article}
\usepackage[margin=1in]{geometry}
\usepackage{amsmath,amssymb,amsthm}
\usepackage{hyperref}

\newtheorem{theorem}{Theorem}
\newtheorem{definition}{Definition}
\newtheorem{example}{Example}

\title{Independence}
\author{math-foundations / node 21}
\date{}

\begin{document}
\maketitle

\section{Definitions}

Let $(\Omega, \mathcal{F}, P)$ be a probability space.

\begin{definition}[Pairwise independence]
Events $A, B \in \mathcal{F}$ are \emph{independent} iff
\[
P(A \cap B) = P(A)\, P(B).
\]
\end{definition}

\begin{definition}[Mutual independence]
A finite family $\{A_1, \dots, A_n\} \subset \mathcal{F}$ is \emph{mutually independent} iff for every
non-empty subset $S \subseteq \{1, \dots, n\}$,
\[
P\!\left( \bigcap_{i \in S} A_i \right) = \prod_{i \in S} P(A_i).
\]
Mutual independence is strictly stronger than pairwise independence.
\end{definition}

\begin{definition}[Conditional independence]
Given $C \in \mathcal{F}$ with $P(C) > 0$, events $A, B$ are \emph{conditionally independent given} $C$ iff
\[
P(A \cap B \mid C) = P(A \mid C)\, P(B \mid C).
\]
\end{definition}

\section{Main theorem}

\begin{theorem}[Independence of complements]
If $A$ and $B$ are independent, then so are $A^c$ and $B$.
\end{theorem}

\begin{proof}
Decompose $B$ as a disjoint union $B = (A \cap B) \sqcup (A^c \cap B)$. By finite additivity,
\[
P(B) = P(A \cap B) + P(A^c \cap B),
\]
so
\[
P(A^c \cap B) = P(B) - P(A \cap B).
\]
Using independence of $A$ and $B$,
\[
P(A^c \cap B) = P(B) - P(A)\, P(B) = (1 - P(A))\, P(B) = P(A^c)\, P(B).
\]
Hence $A^c$ and $B$ are independent.
\end{proof}

Applying the theorem twice (first to $(A,B)$, then to $(B, A^c)$) shows that $A^c$ and $B^c$ are also independent.

\section{A pairwise-but-not-mutually-independent triple}

\begin{example}
Let $\Omega = \{1,2,3,4\}$ with the uniform measure. Define
\[
A = \{1,2\}, \quad B = \{1,3\}, \quad C = \{1,4\}.
\]
Then $P(A) = P(B) = P(C) = \tfrac{1}{2}$. Each pairwise intersection equals $\{1\}$, so
$P(A \cap B) = P(A \cap C) = P(B \cap C) = \tfrac{1}{4} = \tfrac{1}{2} \cdot \tfrac{1}{2}$, giving pairwise independence.
But $A \cap B \cap C = \{1\}$, so $P(A \cap B \cap C) = \tfrac{1}{4} \neq \tfrac{1}{8} = P(A) P(B) P(C)$.
\end{example}

\end{document}
